% =====================================================================
% Exercice 4 — Router LCEL avec RunnableBranch et classifieur LLM
% Compilation : pdflatex -shell-escape exercice4-router.tex
% (le -shell-escape est requis par minted pour la coloration syntaxique)
% Pré-requis : pygmentize (pip install Pygments) + paquet TeX "minted"
% =====================================================================
\documentclass[11pt,a4paper]{article}

% ---------- Encodage & langue ----------
\usepackage[utf8]{inputenc}
\usepackage[T1]{fontenc}
\usepackage[french]{babel}
\usepackage{lmodern}
\usepackage{microtype}

% ---------- Mise en page ----------
\usepackage[a4paper,margin=2.2cm]{geometry}
\usepackage{parskip}
\usepackage{xcolor}
\usepackage{enumitem}
\usepackage{hyperref}
\hypersetup{
  colorlinks=true,
  linkcolor=blue!60!black,
  urlcolor=blue!60!black,
  citecolor=blue!60!black,
}

% ---------- Boîtes & callouts ----------
\usepackage[most]{tcolorbox}
\tcbset{
  enhanced,
  colback=blue!3,
  colframe=blue!50!black,
  fonttitle=\bfseries,
  arc=2pt, boxrule=0.6pt,
}
\newtcolorbox{idee}[1][]{title=Idée clé,colback=yellow!8,colframe=orange!70!black,#1}
\newtcolorbox{piege}[1][]{title=Piège,colback=red!5,colframe=red!60!black,#1}
\newtcolorbox{wow}[1][]{title=Le « wow » moment,colback=green!5,colframe=green!50!black,#1}

% ---------- Code avec minted ----------
\usepackage{minted}
\setminted{
  fontsize=\small,
  linenos=true,
  numbersep=6pt,
  frame=lines,
  framesep=2mm,
  bgcolor=gray!4,
  breaklines=true,
  breakanywhere=true,
  python3=true,
  tabsize=4,
}
\newminted{python}{}

% ---------- TikZ ----------
\usepackage{tikz}
\usetikzlibrary{
  arrows.meta,
  positioning,
  shapes.geometric,
  shapes.misc,
  fit,
  backgrounds,
  calc,
  trees,
}

% Styles TikZ communs
\tikzset{
  node distance=8mm and 12mm,
  block/.style={
    rectangle, rounded corners=2pt, draw=blue!60!black, fill=blue!8,
    minimum width=42mm, minimum height=9mm, align=center, font=\small
  },
  data/.style={
    rectangle, rounded corners=1pt, draw=gray!60, fill=gray!10,
    minimum width=42mm, minimum height=8mm, align=center, font=\footnotesize\ttfamily
  },
  predicate/.style={
    diamond, draw=orange!70!black, fill=orange!15, aspect=2.4,
    inner sep=1pt, font=\footnotesize, align=center
  },
  llm/.style={
    rectangle, rounded corners=2pt, draw=red!60!black, fill=red!8,
    minimum width=42mm, minimum height=9mm, align=center, font=\small
  },
  flow/.style={-{Stealth[length=2.2mm]}, thick, draw=blue!60!black},
  flowdata/.style={-{Stealth[length=2.2mm]}, thick, draw=gray!60, dashed},
  label/.style={font=\scriptsize\itshape, fill=white, inner sep=1pt},
}

% ---------- Titre ----------
\title{\vspace{-1.5em}\textbf{Exercice 4 — Router LCEL}\\
\large Aiguillage dynamique avec \texttt{RunnableBranch} et un classifieur LLM}
\author{Cours PyTutor2 — Phase 1 (LangChain LCEL)}
\date{}

\begin{document}
\maketitle

% =====================================================================
\section*{Objectifs pédagogiques}
\begin{itemize}[leftmargin=*]
  \item Comprendre le \textbf{pattern \emph{router}}~: faire dispatcher une
        requête vers plusieurs sous-chaînes spécialisées selon son contenu.
  \item Manipuler trois briques fondamentales de LCEL~:
        \texttt{RunnablePassthrough.assign}, \texttt{RunnableBranch} et
        \texttt{with\_structured\_output}.
  \item Vérifier visuellement le comportement obtenu dans \textbf{LangSmith}
        (arbre d'exécution, prompts effectifs, tokens, coûts).
\end{itemize}

% =====================================================================
\section{Énoncé}

On veut un assistant Python qui adapte son ton selon la nature du concept
posé en question. Trois personnalités sont attendues~:

\begin{itemize}[leftmargin=*]
  \item \textbf{syntaxe}~: pour les éléments du langage
        (\texttt{def}, \texttt{lambda}, \texttt{yield}, comprehensions, \dots).
        Doit montrer un mini-exemple et un anti-exemple.
  \item \textbf{paradigme}~: pour les concepts de design (POO, polymorphisme,
        fonctionnel, \dots). Doit insister sur la philosophie et les implications.
  \item \textbf{ecosysteme}~: pour les outils / bibliothèques (\texttt{pip},
        \texttt{venv}, \texttt{pytest}, \texttt{numpy}, \dots). Doit donner
        les commandes shell pratiques.
\end{itemize}

\textbf{Contrainte}~: tout doit rester un seul \texttt{Runnable} LCEL pour
préserver le streaming, le batching, et un tracing LangSmith unifié.

% =====================================================================
\section{Architecture}

\subsection*{Vue d'ensemble}

\begin{center}
\begin{tikzpicture}
  \node[data] (in) {\{"concept": "lambda"\}};
  \node[block, below=of in] (assign) {RunnablePassthrough.assign\\(category=classifier)};
  \node[data, below=of assign] (enriched)
    {\{"concept": "lambda",\\\ "category": "syntaxe"\}};
  \node[block, below=of enriched] (branch) {RunnableBranch\\(3 cas + défaut)};
  \node[llm, below=of branch] (out) {réponse spécialisée\\(AIMessage)};

  \draw[flowdata] (in)       -- (assign);
  \draw[flow]     (assign)   -- (enriched) node[label, midway, right]{enrichit le dict};
  \draw[flowdata] (enriched) -- (branch);
  \draw[flow]     (branch)   -- (out) node[label, midway, right]{choisit la sous-chaîne};
\end{tikzpicture}
\end{center}

L'entrée est un dict avec une seule clé \texttt{concept}. Un \emph{classifieur}
(un premier LLM contraint à 3 valeurs) ajoute la clé \texttt{category}. Le
\texttt{RunnableBranch} regarde cette clé et exécute la sous-chaîne adéquate.

% =====================================================================
\section{Le code complet}

\begin{pythoncode}
class CategoryResult(BaseModel):
    """Catégorie d'un concept Python pour router vers la bonne sous-chaîne."""

    category: Literal["syntaxe", "paradigme", "ecosysteme"] = Field(
        description=(
            "Catégorise le concept :\n"
            " - 'syntaxe'    : éléments du langage (def, lambda, comprehensions, with, yield...)\n"
            " - 'paradigme'  : concepts de design (POO, fonctionnel, héritage, polymorphisme...)\n"
            " - 'ecosysteme' : outils/bibliothèques (pip, venv, pytest, numpy, requests...)"
        )
    )


def build_classifier():
    """Mini-chaîne qui catégorise un concept en 3 classes."""
    classify_prompt = ChatPromptTemplate.from_messages([
        ("system", "Tu classes des concepts Python en 3 catégories. "
                   "Ne réponds que par la catégorie."),
        ("human",  "Concept à catégoriser : {concept}"),
    ])
    # On récupère .category (le string) plutôt que l'objet entier pour pouvoir
    # l'utiliser facilement dans les prédicats de RunnableBranch.
    return (
        classify_prompt
        | get_model().with_structured_output(CategoryResult)
        | RunnableLambda(lambda r: r.category)
    )


def build_specialized_chain(focus: str):
    """Construit une sous-chaîne avec un prompt système orienté `focus`."""
    prompt = ChatPromptTemplate.from_messages([
        ("system",
         "Tu es un assistant Python. Le concept à expliquer relève de la "
         f"catégorie '{focus}'. Adapte ton explication en conséquence :\n" +
         {
             "syntaxe":    "insiste sur la syntaxe exacte, montre un mini exemple et un anti-exemple.",
             "paradigme":  "insiste sur la philosophie, le 'pourquoi', les implications de design.",
             "ecosysteme": "insiste sur l'usage pratique, l'installation, les commandes shell typiques.",
         }[focus]),
        ("human", "Explique : {concept}"),
    ])
    return prompt | get_model()


def build_router():
    """Le routeur complet : classifie puis branche."""
    classifier        = build_classifier()
    syntaxe_chain     = build_specialized_chain("syntaxe")
    paradigme_chain   = build_specialized_chain("paradigme")
    ecosysteme_chain  = build_specialized_chain("ecosysteme")

    branch = RunnableBranch(
        (lambda x: x["category"] == "syntaxe",   syntaxe_chain),
        (lambda x: x["category"] == "paradigme", paradigme_chain),
        ecosysteme_chain,  # défaut : si rien ne matche
    )

    # On enrichit le dict d'entrée avec la catégorie, puis on branche.
    return RunnablePassthrough.assign(category=classifier) | branch
\end{pythoncode}

% =====================================================================
\section{Explications pas à pas}

\subsection{Étape 0 — Pourquoi ce pattern existe}

Un seul prompt système ne peut pas être bon pour tout~: \emph{« explique
\texttt{lambda} »} mérite un focus sur la syntaxe ; \emph{« explique le
polymorphisme »} mérite une réflexion conceptuelle ; \emph{« explique pytest »}
mérite des commandes shell. Plutôt qu'un méga-prompt qui essaie de tout couvrir,
on \textbf{spécialise} puis on \textbf{route}.

\begin{idee}
Cet exercice fait \textbf{2 appels LLM} pour répondre à une question~:
un premier pour décider la catégorie, un second (au prompt spécialisé)
pour répondre. C'est le surcoût à accepter pour gagner en qualité de réponse.
\end{idee}

\subsection{Étape 1 — Le vocabulaire de routage (\texttt{CategoryResult})}

Avant de router, il faut un \textbf{ensemble fermé} de catégories. Le
\texttt{Literal["syntaxe", "paradigme", "ecosysteme"]} est crucial~: c'est lui
qui garantit que le LLM ne pourra \emph{physiquement pas} renvoyer
\texttt{"truc"} ou \texttt{"je sais pas"}. La \texttt{description} du champ
sert de mini-prompt qui explique au modèle \emph{comment} classer.

\begin{idee}
Avec \texttt{with\_structured\_output}, le schéma Pydantic \textbf{est} une
partie du prompt. Soigne les \texttt{description=} des champs autant que tu
soignerais un message système.
\end{idee}

\subsection{Étape 2 — Le classifieur (\texttt{build\_classifier})}

\begin{center}
\begin{tikzpicture}[node distance=6mm and 14mm]
  \node[data] (in) {\{"concept": "lambda"\}};
  \node[block, right=of in] (prompt) {ChatPromptTemplate};
  \node[llm,   right=of prompt] (llm) {ChatAnthropic\\with\_structured\_output};
  \node[block, below=of llm] (unwrap) {RunnableLambda\\(lambda r: r.category)};
  \node[data,  left=of unwrap] (out) {"syntaxe"};

  \draw[flowdata] (in)     -- (prompt);
  \draw[flow]     (prompt) -- (llm)
    node[label, midway, above, sloped]{messages};
  \draw[flow]     (llm)    -- (unwrap)
    node[label, midway, right]{CategoryResult(category="syntaxe")};
  \draw[flow]     (unwrap) -- (out)
    node[label, midway, above]{string};
\end{tikzpicture}
\end{center}

Trois maillons~:

\begin{enumerate}[leftmargin=*]
  \item \texttt{classify\_prompt}~: \texttt{\{"concept": "lambda"\}}
        $\rightarrow$ messages formatés.
  \item \texttt{model.with\_structured\_output(CategoryResult)}~: messages
        $\rightarrow$ objet \texttt{CategoryResult(category="syntaxe")}.
  \item \texttt{RunnableLambda(lambda r: r.category)}~: objet $\rightarrow$
        \texttt{"syntaxe"} (string).
\end{enumerate}

Le maillon 3 n'est pas cosmétique~: il rend la sortie utilisable comme
\textbf{valeur de comparaison} dans les prédicats du \texttt{RunnableBranch}.

\subsection{Étape 3 — Les sous-chaînes spécialisées}

\texttt{build\_specialized\_chain(focus)} est une \emph{fabrique}~: on
l'appelle 3 fois avec des \texttt{focus} différents et on obtient 3 chaînes
distinctes, chacune avec sa propre personnalité dans le system prompt.

\begin{idee}
La spécialisation se fait à la \textbf{construction}, pas à l'appel.
\texttt{build\_specialized\_chain("syntaxe")} renvoie un \texttt{Runnable}
prêt à l'emploi qui ne sait répondre que dans le style « syntaxe ».
\end{idee}

\subsection{Étape 4 — L'assemblage (\texttt{build\_router})}

\subsubsection*{4a — La brique de branchement}

\texttt{RunnableBranch} est le \texttt{if/elif/else} de LCEL~:

\begin{center}
\begin{tikzpicture}[node distance=5mm and 12mm]
  \node[data] (in) {\{"concept": "lambda",\\\ "category": "syntaxe"\}};

  \node[predicate, right=18mm of in] (p1) {category\\=="syntaxe"?};
  \node[predicate, below=10mm of p1] (p2) {category\\=="paradigme"?};
  \node[block,     below=10mm of p2] (def) {ecosysteme\_chain\\(défaut)};

  \node[llm, right=of p1] (c1) {syntaxe\_chain};
  \node[llm, right=of p2] (c2) {paradigme\_chain};

  \draw[flowdata] (in) -- (p1);
  \draw[flow]     (p1) -- node[label, above]{vrai}   (c1);
  \draw[flow]     (p1) -- node[label, right]{faux}   (p2);
  \draw[flow]     (p2) -- node[label, above]{vrai}   (c2);
  \draw[flow]     (p2) -- node[label, right]{faux}   (def);
\end{tikzpicture}
\end{center}

Points clés~:
\begin{itemize}[leftmargin=*]
  \item Les prédicats sont évalués \textbf{dans l'ordre}, le premier qui
        renvoie \texttt{True} gagne (court-circuit comme un \texttt{if/elif}).
  \item Si \emph{aucun} prédicat ne matche, la chaîne par défaut s'exécute.
        Elle est \textbf{obligatoire}.
  \item La sous-chaîne sélectionnée reçoit \textbf{l'input complet} (le dict
        entier), pas juste le résultat du prédicat. C'est pour ça que
        \texttt{syntaxe\_chain} peut toujours lire \texttt{\{concept\}}
        dans son prompt.
\end{itemize}

\subsubsection*{4b — L'enrichissement avant branchement}

Le \texttt{RunnableBranch} a besoin que \texttt{category} \textbf{soit déjà
dans l'input} pour que ses prédicats puissent le lire. C'est
\texttt{RunnablePassthrough.assign} qui s'en charge~:

\begin{center}
\begin{tikzpicture}[node distance=6mm and 0mm]
  \node[data] (in) {\{"concept": "lambda"\}};
  \node[block, below=of in] (assign)
    {RunnablePassthrough.assign(category=classifier)};
  \node[data, below=of assign] (mid)
    {\{"concept": "lambda", "category": "syntaxe"\}};
  \node[block, below=of mid] (branch) {RunnableBranch};
  \node[llm, below=of branch] (out) {AIMessage("Une lambda est\dots")};

  \draw[flowdata] (in) -- (assign);
  \draw[flow]     (assign) -- (mid)
    node[label, midway, right]{appel classifier inclus ici};
  \draw[flowdata] (mid) -- (branch);
  \draw[flow]     (branch) -- (out)
    node[label, midway, right]{choisit syntaxe\_chain};
\end{tikzpicture}
\end{center}

\begin{idee}
\texttt{RunnablePassthrough.assign} est la manière idiomatique LCEL
d'\textbf{enrichir} un dict d'étape en étape, comme un
\texttt{SELECT *, expr AS col} en SQL~: les clés existantes sont conservées,
on ajoute juste les nouvelles.
\end{idee}

% =====================================================================
\section{Exécution et observation}

\subsection*{Le code de l'exécuteur (\texttt{exercise\_4})}

\begin{pythoncode}
def exercise_4() -> None:
    router = build_router()

    test_concepts = [
        "les list comprehensions",   # syntaxe attendue
        "le polymorphisme",          # paradigme attendu
        "pytest",                    # écosystème attendu
        "lambda",                    # syntaxe
    ]

    # Pour la démo, on veut voir la catégorie choisie. On reconstruit un
    # pipeline qui garde la catégorie dans la sortie.
    classifier = build_classifier()
    debug_chain = RunnablePassthrough.assign(
        category=classifier,
        answer=router,
    )

    for concept in test_concepts:
        result = debug_chain.invoke({"concept": concept})
        category = result["category"]
        snippet = result["answer"].content[:200].replace("\n", " ")
        print(f"\n[{category:>10}] {concept}")
        print(f"             → {snippet}...")
\end{pythoncode}

\begin{piege}
Ce \texttt{debug\_chain} appelle le classifieur \textbf{deux fois} par concept
(une fois pour la clé \texttt{category} directe, une fois enfoui dans
\texttt{router}). Donc \textbf{3 appels LLM par concept} au total, au lieu de 2.
C'est acceptable pour la démo ; en production on factorise.
\end{piege}

% =====================================================================
\section{Vérification dans LangSmith}

\subsection{Configuration minimale}

Ajouter dans \texttt{.env}~:

\begin{minted}{bash}
LANGSMITH_TRACING=true
LANGSMITH_API_KEY=lsv2_pt_...
LANGSMITH_PROJECT=PyTutor2
# LANGSMITH_ENDPOINT=https://eu.api.smith.langchain.com  # uniquement compte EU
\end{minted}

Puis~:
\begin{minted}{bash}
python phase1-solutions.py
\end{minted}

Aller sur \url{https://smith.langchain.com} $\rightarrow$ projet
\texttt{PyTutor2}. Quatre runs apparaissent (un par concept de la liste).

\subsection{L'arbre d'exécution attendu}

Pour le concept \texttt{"lambda"}, le panneau gauche affiche~:

\begin{center}
\begin{tikzpicture}[
  grow via three points={one child at (0.5,-0.55) and two children at (0.5,-0.55) and (0.5,-1.1)},
  edge from parent path={(\tikzparentnode.south) |- (\tikzchildnode.west)},
  every node/.style={font=\footnotesize\ttfamily, anchor=west},
  level distance=0pt,
  sibling distance=0pt,
]
\node {RunnableSequence  (debug\_chain, $\sim$2{-}4 s, $\sim$1500 tokens)}
  child {node {RunnableAssign<category, answer>}
    child {node {RunnableParallel<category, answer>}
      child {node {[category] RunnableSequence}
        child {node {ChatPromptTemplate}}
        child {node {ChatAnthropic (structured\_output)  ← appel LLM \#1}}
        child {node {RunnableLambda  → "syntaxe"}}
      }
      child [missing] {}
      child [missing] {}
      child [missing] {}
      child {node {[answer] RunnableSequence (le router)}
        child {node {RunnableAssign<category>}
          child {node {(ré-exécution du classifier) ← appel LLM \#2}}
        }
        child [missing] {}
        child {node {RunnableBranch}
          child {node {branch[0]: syntaxe\_chain (choisie)}
            child {node {ChatPromptTemplate}}
            child {node {ChatAnthropic  ← appel LLM \#3}}
          }
        }
      }
    }
  };
\end{tikzpicture}
\end{center}

\subsection{Les quatre choses à regarder}

\begin{enumerate}[leftmargin=*]
  \item \textbf{L'arbre lui-même}~: tu \emph{vois} 3 appels LLM, pas 2. C'est
        exactement la fuite de coût que je mentionne dans la mise en garde
        ci-dessus.
  \item \textbf{Premier \texttt{ChatAnthropic}} (classifieur)~: onglet
        \emph{Input} → messages formatés ; onglet \emph{Output} → l'objet
        \texttt{CategoryResult} reconstitué, et le tool-call sous-jacent que
        Claude a utilisé pour respecter le schéma.
  \item \textbf{Deuxième \texttt{ChatAnthropic}} (réponse spécialisée)~: le
        system prompt commence par \emph{« Tu es un assistant Python\dots
        catégorie 'syntaxe' »}. Tu vois \textbf{le prompt exact} qui a produit
        la réponse — précieux pour débugger.
  \item \textbf{Onglet \emph{Costs}} sur le run racine~: total tokens et coût
        \$ agrégés sur les 3 appels. Permet de chiffrer le surcoût du pattern
        router par rapport à un single-prompt.
\end{enumerate}

\begin{wow}
Compare deux runs dans l'UI~:
\begin{itemize}[leftmargin=*]
  \item \texttt{"lambda"} $\rightarrow$ la branche \texttt{syntaxe\_chain} s'exécute,
  \item \texttt{"pytest"} $\rightarrow$ la branche \texttt{ecosysteme\_chain} s'exécute.
\end{itemize}
Dans \texttt{RunnableBranch}, tu vois littéralement \textbf{une branche choisie}
selon le concept. C'est la visualisation directe de l'aiguillage dynamique.
\end{wow}

% =====================================================================
\section{Pour aller plus loin}

\subsection*{Optimisation~: supprimer le double appel}

Le \texttt{classifier} tourne deux fois par concept à cause du \texttt{debug\_chain}.
On peut le calculer une seule fois et \emph{partager} le résultat~:

\begin{pythoncode}
def build_router_observable():
    classifier = build_classifier()
    syntaxe_chain    = build_specialized_chain("syntaxe")
    paradigme_chain  = build_specialized_chain("paradigme")
    ecosysteme_chain = build_specialized_chain("ecosysteme")

    branch = RunnableBranch(
        (lambda x: x["category"] == "syntaxe",   syntaxe_chain),
        (lambda x: x["category"] == "paradigme", paradigme_chain),
        ecosysteme_chain,
    )

    # Un seul .assign(category=...) : la clé est ensuite disponible pour
    # le branch ET pour la sortie finale.
    return (
        RunnablePassthrough.assign(category=classifier)
        | RunnablePassthrough.assign(answer=branch)
    )
\end{pythoncode}

Ici on passe à \textbf{2 appels LLM par concept} (au lieu de 3), tout en
gardant \texttt{category} dans la sortie finale.

\subsection*{Exercices proposés}

\begin{enumerate}[leftmargin=*]
  \item Ajouter une 4\textsuperscript{e} catégorie \texttt{"performance"} (ex.~:
        \emph{GIL}, \emph{multiprocessing}, \emph{asyncio}). Adapter le
        \texttt{Literal}, ajouter une sous-chaîne, ajouter un prédicat dans
        le branch.
  \item Détecter dans LangSmith les concepts qui ont été \emph{mal classés}.
        Annoter manuellement quelques runs avec un \emph{feedback} et constater
        la trace dans l'onglet \emph{Annotations}.
  \item Mesurer le coût moyen par concept via l'API LangSmith
        (\texttt{client.list\_runs(...)}) et calculer le surcoût du router vs
        un single-prompt « universel ».
\end{enumerate}

% =====================================================================
\section*{Récapitulatif en une page}

\begin{tcolorbox}[title=Ce qu'il faut retenir]
\begin{enumerate}[leftmargin=*]
  \item \textbf{Pattern router} = classifieur LLM + sous-chaînes spécialisées
        + \texttt{RunnableBranch}.
  \item \textbf{\texttt{with\_structured\_output} + \texttt{Literal}} =
        garantie que la sortie est dans un ensemble fermé, donc routable.
  \item \textbf{\texttt{RunnablePassthrough.assign}} = manière LCEL d'enrichir
        le dict de contexte au fil de la pipeline, sans perdre les clés
        précédentes.
  \item \textbf{Tout reste un seul \texttt{Runnable}}~: \texttt{router.invoke},
        \texttt{router.batch}, \texttt{router.astream} marchent — et LangSmith
        trace l'arbre complet en un seul run.
\end{enumerate}
\end{tcolorbox}

\end{document}